\section{Proofs of the worked examples}\label{app:examples}

This appendix derives the evaluations in Section~\ref{sec:examples}
from the saddle and the named endpoint definitions.  We first record the
two observations that permit exact face calculations and the passage
from positive vectors to arithmetic families with repeated cutoffs.

\subsection{Stationarity and bounded coordinates}

Suppose a candidate saddle has a proper owner \([1,n]\) with product
\(s\), followed by a singleton terminal owner with product \(\theta<s\).
Write \(z=2^{\theta/s}\), \(r_i=n-i+1\), and
\(\kappa_r=(r+z)^{s-1}\).  On this face the two common-direction
stationarity equations are
\begin{align}
 &\sum_{i=1}^n L_i\{\kappa_{r_i}E_{r_i}(z)-r_i\}=0,
 \label{exapp:owner-balance}\\
 &L_{n+1}(2^\theta\log2-1)
   =\sum_{i=1}^nL_i\{1-z\log2\,\kappa_{r_i}\}.
 \label{exapp:child-balance}
\end{align}
For \(1\le j<n\), the internal multiplier is
\begin{equation}
 \lambda_j=
 \sum_{i=1}^j L_i
 \{\kappa_{r_i}[E_{r_i}(z)-E_{n-j}(z)]-(r_i-n+j)\}.
 \label{exapp:prefix-KKT}
\end{equation}
These formulas follow by differentiating \(G_i=s\log(r_i+z)\).
When differentiating in \(s\), the child product \(\theta\) is held
fixed, so the derivative of \(z\) must be included.  For example,
\[
 \frac{\partial}{\partial s}(r+2^{\theta/s})^s
       =(r+z)^{s-1}E_r(z).
\]
For the prefix in \eqref{exapp:prefix-KKT}, the child has weight
\(n-j+z\); taking the difference of the corresponding entropy
increments gives the displayed formula.  Its sign convention agrees
with \eqref{hier:eq:KKT}.  Thus nonnegative internal multipliers,
zero exterior multipliers, and strict separation of the two products
verify all KKT conditions.  Theorem~\ref{geom:thm:global} then identifies
the candidate as the unique global saddle.  The same calculation for a
single terminal owner uses \(z=1\), with common-direction balance
\begin{equation}
 \sum_iL_i\{h_i^s\log h_i-(h_i-1)\}=0.
 \label{exapp:terminal-balance}
\end{equation}

We use the word \emph{saturated} below in the asymptotic sense that a
binary factor has a positive lower bound independent of the parameter
tending to infinity.  It need not equal one.  If one endpoint score is
\(\gg L_i\), the other is at least one, and the native intensity remains
in a fixed compact interval, the factor for group \(i\) is saturated.
Every edge on a bounded positive clock is saturated for the same reason.

Choose \(C\ge L_*(k)\), where \(L_*(k)\) is the fixed threshold in
Theorem~\ref{lower:lifted}.  On vectors all of whose coordinates exceed
this threshold, that theorem and Theorem~\ref{positive:upper} give
\(K_k\asymp_k\Psi_k\); Theorem~\ref{main:uniform} then gives
\(\Phi_k\asymp_k\Psi_k\).  Replacing a fixed number of zero coordinates
by \(C\) changes the total log-parameter distance by a bounded amount.
Theorem~\ref{thm:ell-stability} therefore transfers the arithmetic
estimate to the original vector.  This proves the required statements
about \(\Phi_k\) without any assertion about where its defining
infimum is attained.  In particular, a raw saddle is never assigned
to a vector with zero coordinates.

\subsection{A finite rational calculation used below}
\label{exapp:certificate-method}

All decimal endpoints in this appendix denote exact rational numbers.
We give the finite inequalities explicitly and describe a single
rational procedure for verifying them.  For \(1\le x\le2\), let
\(w=(x-1)/(x+1)\).  Then
\begin{equation}
 2\sum_{j=0}^{N-1}\frac{w^{2j+1}}{2j+1}
 \le\log x\le
 2\sum_{j=0}^{N-1}\frac{w^{2j+1}}{2j+1}
 +\frac{2w^{2N+1}}{(2N+1)(1-w^2)}.
 \label{exapp:log-bounds}
\end{equation}
This follows by integrating the geometric series for
\(2/(1-w^2)\), and bounding all remaining denominators below by
\(2N+1\).  General positive rational arguments are reduced to
\([1,2]\) by extracting powers of two and, if necessary, taking
reciprocals.  For \(0\le x<N+2\), Taylor's series gives
\begin{equation}
 \sum_{j=0}^N\frac{x^j}{j!}\le e^x\le
 \sum_{j=0}^N\frac{x^j}{j!}
 +\frac{x^{N+1}}{(N+1)!}\frac{1}{1-x/(N+2)}.
 \label{exapp:exp-bounds}
\end{equation}
Indeed, successive omitted terms have ratio at most \(x/(N+2)\).
Reduce a nonnegative argument to \([0,1/4]\), then square back;
negative arguments are handled by reciprocation.

The bounds below result from \(48\) terms in
\eqref{exapp:log-bounds}, degree \(40\) in
\eqref{exapp:exp-bounds}, and outward rounding to the rational grid
\(10^{-40}\mathbb Z\) after each operation.  Interval addition,
multiplication, and division use their elementary rational endpoint
formulas; divisions are used only when zero is excluded from the
denominator interval.  Monotonicity supplies interval logarithms and
exponentials.  Thus each table is a finite rational inequality, with
explicit remainders.  Exact zeros used in the constructions are instead
proved algebraically from their defining equations.

The finite sums and interval operations specified here suffice to
reproduce every bound below.  Neither numerical optimization nor
floating-point equality tests are needed.

\subsection{Equal cutoffs}

Evaluate the positive vector \(Q=(T,C,\ldots,C)\).  At a common
product \(s\), equation \eqref{exapp:terminal-balance} is strictly
increasing in \(s\), and gives
\(s=\alpha_H+O_k(T^{-1})\).  For \(1\le j<k\), its multiplier is
\begin{equation}
 \lambda_j=T\left[
 H^{\alpha_H-1}\{H\log H-(H-j)\log(H-j)\}-j\right]+O_k(1).
 \label{exapp:equal-KKT}
\end{equation}
The term in brackets is positive.  To prove this, apply strict
convexity of \(x\log x\): the secant over \([H-j,H]\) has larger
slope than the secant over \([1,H]\).  The latter equals
\(H^{1-\alpha_H}\), since
\(H^{\alpha_H}\log H=H-1\).  Multiplication by
\(jH^{\alpha_H-1}\) gives the claim.  All internal multipliers are
therefore positive for large \(T\), and common-direction stationarity
proves the asserted single owner.

For each first-group binary edge \(k\to q\), \(1\le q<k\), the left
completion retains exactly \(q\) names through the first group.  Its
reference cost is
\[
 T\{q-H^{\alpha_H-1}(q+1)\log(q+1)\}+O_k(1)\gg_k T.
\]
The inequality follows because \((q+1)\log(q+1)/q\) is strictly
increasing and the corresponding quantity at \(q=k\) is
\(H^{1-\alpha_H}\).  The other clocks are fixed, so all binary
factors are saturated.  The ordinary suffix starts in group \(1\),
because \(|s-\alpha_H|<\tau\) for large \(T\).  Equation
\eqref{exapp:terminal-balance} gives \(D=0\), while
\(U=T+(k-1)C\), \(M=T\), \(A=0\), and \(B=(k-1)C\).
Formula \eqref{profile:ordinary} gives \(T^{-3/2}\).
Finally \(J(Q)=I_HT+O_k(1)\).  The bounded-coordinate argument above
proves \eqref{ex:equal-answer}.

\subsection{The critical family with six divisor coordinates}

Use \(\sigma_c,e_h,a,b\) from
\eqref{ex:critical-constants}--\eqref{ex:critical-lengths}.  Put
\[
 f_h(q)=q-h^{\sigma_c-1}(q+1)\log(q+1),\qquad
 \Delta_h(Q)=h^{\sigma_c-1}(h\log h-Q\log Q)-(h-Q).
\]
Thus \(f_h(0)=0\), \(f_h(h-1)=-e_h\), and
\(\Delta_h(Q)=f_h(Q-1)-f_h(h-1)\).  The finite inequalities needed
here are
\begin{gather}
 e_2>0,\qquad e_h<0\quad(3\le h\le7),\nonumber\\
 \Delta_7(2)=0,\qquad
 \Delta_7(Q)>0\quad(3\le Q\le6),\qquad
 \Delta_h(Q)>0\quad(3\le h\le6,\ 2\le Q<h).
 \label{exapp:critical-signs}
\end{gather}
The zero is exactly the definition of \(\sigma_c\).  The rational
bounds of \S\ref{exapp:certificate-method} give
\[
 \begin{array}{c|c}
 \text{quantity}&\text{enclosure or lower bound}\\ \hline
 \sigma_c&(0.540128277594,\ 0.540128277595)\\
 e_2&(0.00790656,\ 0.00790657)\\
 a&(54.82474304,\ 54.82474305)\\
 b&(66.42545992,\ 66.42545993)\\
 \min\{\Delta_h(Q):(h,Q)\ne(7,2)\text{ in \eqref{exapp:critical-signs}}\}
       &>0.1079
 \end{array}
\]
They also give
\(\sigma_c-\alpha_2>1/(8\cdot7^3(\log7)^2)\) and
\(\alpha_3-\sigma_c>1/(8\cdot7^3(\log7)^2)\).
Since the right side is an upper bound for the profile's \(\tau\),
the ordinary suffix is exactly the last group.  The signs of \(e_h\)
also follow from these two inequalities and the strict increase of
\(\alpha_h\).

The common-direction balance holds exactly by the definitions of
\(a,b\).  For \(1\le j\le5\), the prefix derivative gives
\[
 \lambda_j=\sum_{i=1}^jL_i\Delta_{8-i}(7-j).
\]
For \(j\le4\), each summand is positive by
\eqref{exapp:critical-signs}.  At \(j=5\), the first summand is zero
and
\(\lambda_5=C\sum_{h=3}^6\Delta_h(2)>0\).
Together with \(\lambda_6=0\), these verify the full KKT system.
In particular the fixed positive buffers keep the last internal
multiplier positive; the actual family is not an owner-merging boundary.

We next verify both endpoint minima, rather than just exhibit one
favorable pair.  For the edge \(6\to1\), the only freely chosen cut
is \(\tau_6\in\{V_1,\ldots,V_6\}\).  If \(\tau_6<V_6\), there is
no negative contribution from the final group.  The function \(f_h\)
is strictly concave, vanishes at zero, and is positive at \(h-1\)
for \(h\ge3\).  It is therefore positive at every intermediate
integer.  The first-group contribution is \(\gg T\) on both sides.
If \(\tau_6=V_6\), use total balance to subtract the full surviving
configuration.  Since \(\Delta_7(2)=0\), both reference costs are
exactly \(C\sum_{h=3}^6\Delta_h(2)=C\Delta\).  The completion is
late, and its good groups are exactly those listed before
\eqref{ex:critical-endpoints}.  Hence the two scores in that equation
are the separate minima for all sufficiently large \(T\).

For the other edges, a completion with \(\tau_6<V_6\) again has cost
\(\gg T\): a retained name survives the first group, and no negative
last-group contribution occurs.  If \(\tau_6=V_6\), at least one
name survives every group, and total balance gives
\begin{equation}
 \overline C_v=\sum_{i=1}^5L_i
          \Delta_{8-i}(q_{v,i}+1).
 \label{exapp:critical-context}
\end{equation}
Every summand is nonnegative.  In group \(1\), an edge with
\(q\ge2\) has left cost \(\gg T\).  In a middle group, a proper
subset of the names survives on the left, giving left cost \(\gg C\).
These bounds saturate every remaining edge.  The ordinary factor has
\(U=M=L_6\), \(A=B=0\), and \(D=e_2L_6\), and is therefore
\(\asymp T^{-1/2}\).  Formula \eqref{profile:edge-factor} applied
to \eqref{ex:critical-endpoints} supplies \(T^{-\chi_c}\).
This proves \eqref{ex:critical-answer} and the stated explicit
expression for \(J(L)\).

For completeness, let
\(\sigma_c(H)=1-\log((H\log H-2\log2)/(H-2))/\log H\).
The values of \(\sigma_c(H)-\alpha_2\), for \(H=3,4,5,6\), lie
respectively in
\[
 \begin{array}{cc}
 (-0.117568,-0.117566),&(-0.056865,-0.056864),\\
 (-0.024370,-0.024368),&(-0.003499,-0.003498).
 \end{array}
\]
Thus \(H=7\) is the first admissible integer in this critical
first/last-group construction.

\subsection{The separated family with six divisor coordinates}

It is useful to check the dimensional boundary for the same family
with \(n=k-1\) and \(H=n+2\).  Let \(\alpha=\alpha_2\) and define
\[
 f(H)=H^{\alpha-1}(H\log H-2\log2)-(H-2).
\]
The rational calculation gives
\begin{equation}
 0.025<f(6)<0.026,\qquad -0.110<f(7)<-0.109,
 \qquad 2\alpha-1<\alpha(1-\alpha)\log2.
 \label{exapp:split-signs}
\end{equation}
For real \(x\ge2\), direct differentiation gives
\[
 f''(x)=x^{\alpha-2}\left[
 \alpha(\alpha-1)\log x+2\alpha-1
 -\frac{2\log2(1-\alpha)(2-\alpha)}x\right]<0.
\]
Indeed, the first two terms already have negative sum by
\eqref{exapp:split-signs}.  Since \(f(2)=0\), strict concavity and
the two signs imply \(f(H)>0\) for integer \(3\le H\le6\) and
\(f(H)<0\) for integer \(H\ge7\).

At fixed \(\theta=\alpha\), the derivative of
\((n+2^{\theta/s})^s\) is
\((n+z)^{s-1}E_n(z)\), and its second derivative is strictly positive.
For instance, with \(c=\theta\log2\), \(p=z/(n+z)\), and
\(A=\log(n+z)-p\log z\), it equals
\[
 (n+z)^s\{A^2+c^2p(1-p)/s^3\}>0.
\]
Thus the equation \((n+z)^{s-1}E_n(z)=n\) has at most one root.
Its left side minus \(n\) equals \(f(H)\) at \(s=\alpha\)
and is positive at \(s=1\), since
\(E_n(z)=\int_0^n(1+\log(z+t))\,dt>n\).
For \(H\ge7\) this proves existence of a unique root above
\(\alpha\).  For \(k=6\) the interval
\((0.54070,0.54072)\) isolates it by opposite signs at its endpoints.
For \(H\le6\), the last multiplier on the terminal diagonal has
positive limiting coefficient \(f(H)\); the other prefix coefficients
are positive by concavity as well.  This proves the dimensional
assertion for this family.

Now take \(Q=(T,C,C,C,C,B)\), \(B=T^3\), with \(k=6\).
Solve \eqref{exapp:owner-balance}--\eqref{exapp:child-balance} with
\(n=5\).  After division by \(T\) and \(B\), respectively, the
limiting equations are \eqref{ex:split-root} and
\(2^\theta\log2=1\).  Their Jacobian is triangular with positive
diagonal entries.  The implicit function theorem gives
\begin{equation}
 s=s_*+O(T^{-1}+T/B),\qquad
 \theta=\alpha+O(T/B).
 \label{exapp:split-products}
\end{equation}
The exterior multiplier is zero by the first equation.  For
\(1\le j<5\), the leading internal multiplier is
\[
 T\left[(5+z_*)^{s_*-1}
 \{(5+z_*)\log(5+z_*)-(5+z_*-j)\log(5+z_*-j)\}-j\right].
\]
Strict convexity of \(x\log x\), compared with its secant on
\([z_*,5+z_*]\), makes this positive.  The remaining terms are
\(O(1)\).  The products have a fixed limiting gap, so these are
the full KKT conditions for precisely the two asserted owners.

The first edge \(5\to0\) has no outside cuts.  On the left all five
names end at zero; the completion is early and has bounded score,
since the proper exit gap has a positive limit.  On the right they end
at \(V_1=T\), a late cut for the proper owner.  Its reference cost is
\(T\{5-(5+z)^{s-1}E_5(z)\}=O(1)\), by
\eqref{exapp:owner-balance}, and it has no good group.  Thus both
endpoint scores are \(\asymp1\).
For a first-group edge with \(q\ge1\), the left cost is
\[
 T\{q-(5+z_*)^{s_*-1}E_q(z_*)\}+O(1)\gg T.
\]
This follows from strict convexity of \(E_q(z_*)\) in \(q\), or
equivalently from strict increase of \(E_q(z_*)/q\), with equality
only at rank \(5\).  The middle clocks are bounded.  The terminal
singleton has no binary edge.  Hence the sole unsaturated binary
factor has order \(T^{-2\chi_*}\), and the sole proper count has
order \(T^{-s_*/2}\).  The use of the limiting products in these
powers is valid by \eqref{exapp:split-products} and the local
Lipschitz continuity of \(\chi\).

Finally, \eqref{exapp:child-balance} gives
\[
 D=c_*T+O(1+T^2/B),\qquad
 c_*=1-z_*\log2\,(5+z_*)^{s_*-1}>0.
\]
The rational bounds give \(0.4403<c_*<0.4404\); positivity alone
is needed.  Thus
\(\mathcal B_{\rm ord}\asymp T B^{-3/2}\).
At the limiting products the objective is
\(I_2B+I_{\rm p}T+O(1)\).  Taylor expansion at those products has
linear error \(O(1)\) and quadratic error
\(O(T\lvert s-s_*\rvert^2+B\lvert\theta-\alpha\rvert^2)=O(1)\).
The mixed derivatives are \(O(T)\), so their contribution is also
bounded.  This proves \(J(Q)=I_2B+I_{\rm p}T+O(1)\) and
\eqref{ex:split-answer}, including the transfer from \(Q\) to the
original repeated cutoffs.

\subsection{The positive ray}

For \eqref{ex:ray-lengths}, equations
\eqref{exapp:owner-balance}--\eqref{exapp:child-balance} hold by
the definitions of \(a,b\).  Direct use of
\eqref{exapp:prefix-KKT} gives the following enclosures:
\[
 \begin{array}{c|cccc}
 j&1&2&3&4\\ \hline
 \lambda_j/t&(1.77059,1.77060)&(2.97343,2.97344)&
 (3.37910,3.37911)&(2.61456,2.61457)
 \end{array}
\]
The products are strictly separated, so this proves the exact saddle.

Here the complete endpoint verification is small enough to specify
directly.  Represent the cut \(V_d\) of name \(c\) by its integer
index \(d_c\).  For an edge \((i,q)\), choose
\[
 d_c\in\{0,\ldots,c\}\quad(c<i),\qquad
 d_c\in\{i,\ldots,c\}\quad(6-q\le c\le5),
\]
and set \(d_c=i-1\) on the left and \(d_c=i\) on the right for
\(i\le c\le5-q\).  This is exactly the paired-template definition
of the profile, including zero cuts and ties.  With
\(w=(a,1,1,1,1)\), each cost is the explicit expression
\begin{equation}
 \frac{\overline C_v}{t}=
 \sum_{g=1}^5 w_g\left[q_g-(6-g+z)^{s-1}E_{q_g}(z)\right],
 \qquad q_g=\#\{c:d_c\ge g\}.
 \label{exapp:ray-cost}
\end{equation}
The following table gives rational lower bounds for its minimum over
all templates on the indicated side.  The template count makes the
finite set in each row explicit.
\begin{center}
\begin{tabular}{cccrr}
\(i\)&\(q\)&side&templates&lower bound\\ \hline
1&0&right&1&0.5490\\
1&1&left&5&2.4949\\
1&2&left&20&3.3791\\
1&3&left&60&2.9734\\
1&4&left&120&1.7705\\
2&0&left&2&0.5490\\
2&1&left&8&0.6686\\
2&2&left&24&0.4464\\
2&3&left&48&0.1751\\
3&0&left&6&0.4791\\
3&1&left&18&0.4284\\
3&2&left&36&0.1837\\
4&0&left&24&0.3280\\
4&1&left&48&0.1805\\
5&0&left&120&0.1455
\end{tabular}
\end{center}
For example, the first row has all cuts at \(V_1\), so its cost is
\(a[5-(5+z)^{s-1}E_5(z)]\).  Every other row is obtained by the
same explicit sum \eqref{exapp:ray-cost} over the displayed independent
integer choices, with \(\log,\exp\) bounded by
\eqref{exapp:log-bounds}--\eqref{exapp:exp-bounds}.
Thus each edge has one endpoint score \(\gg t\) and the other at
least one, proving saturation of all fifteen factors.
The proper count is \(\asymp t^{-s/2}\).
The terminal group has \(U=M=bt\) and
\(D=(2^\theta\log2-1)bt\asymp t\), so its ordinary factor is
\(\asymp t^{-1/2}\).  Homogeneity gives \(J(L)=tJ_0\).
All lengths tend to infinity, proving \eqref{ex:ray-answer} directly.

\subsection{Two comparisons within one owner}

We first justify all fixed constants in \eqref{ex:same-root}.  Let
\(D_{r,q}(z)=(E_r(z)-E_q(z))/(r-q)\) and
\(t_{r,q}(z)=1-\log D_{r,q}(z)/\log(r+z)\).  Differentiation gives
\begin{equation}
 t_{r,q}'(z)=
 \frac{\log D_{r,q}(z)}{(r+z)\{\log(r+z)\}^2}
 -\frac{\log(r+z)-\log(q+z)}
 {(r-q)D_{r,q}(z)\log(r+z)}.
 \label{exapp:secant-derivative}
\end{equation}
On \(I=[1.98455,1.98457]\), the rational bounds give
\begin{gather*}
 (t_{14,6}-t_{6,0})(1.98455)
       \in(-0.000000199838,-0.000000199837),\\
 (t_{14,6}-t_{6,0})(1.98457)
       \in(0.000000161388,0.000000161389),\\
 0.01806066<(t_{14,6}-t_{6,0})'(z)<0.01806187
       \qquad(z\in I).
\end{gather*}
The intermediate value theorem and strict monotonicity therefore
give exactly one root in \(I\).  The same calculation on \(I\)
gives
\begin{equation}
 \begin{gathered}
 \alpha_2<\theta<\sigma,\qquad
 -0.166<g_1<-0.165,\qquad 1.604<g_{14}<1.605,\\
 \min_{1\le r\le14}(1-z\log2\,\kappa_r)>0.1574.
 \end{gathered}
 \label{exapp:same-fixed-signs}
\end{equation}
The identities \(g_6=0\) and \(f_{14}(6)=g_{14}\) follow exactly
from the two secant equations, since
\((r+z)^{\sigma-1}(E_r-E_q)=r-q\) is equivalent to
\(\sigma=t_{r,q}(z)\).  With
\(\Delta_r=\kappa_r(E_r-E_6)-(r-6)\), further positive bounds are
\[
\begin{array}{c|rrrrrrr}
 r&7&8&9&10&11&12&13\\ \hline
 \Delta_r&
 >0.1725&>0.2765&>0.3250&>0.3277&>0.2916&>0.2222&>0.1238
\end{array}
\]
The two coefficients in \(W\) and the leading coefficient in \(R\)
are enclosed by
\begin{gather*}
 9.6779<-g_{14}/g_1<9.6822,\\
 25.0555<-\frac{\sum_{2\le r\le13,\ r\ne6}g_r}{g_1}<25.0707,\\
 39.1859<\frac{1-z\log2\,\kappa_6}{2^\theta\log2-1}<39.214.
\end{gather*}
Thus all lengths in \eqref{ex:same-lengths} are positive.  Each
\(f_r(q)\) is strictly concave in \(q\).  Its zero at \(q=0\)
and the equalities just proved imply
\begin{equation}
 f_6(q)>0\ (0<q<6),\qquad
 f_{14}(q)>g_{14}\ (6<q<14),\qquad
 f_{14}(q)>0\ (0<q\le14).
 \label{exapp:same-concavity}
\end{equation}

The lengths \(W,R\) give exactly
\eqref{exapp:owner-balance}--\eqref{exapp:child-balance} for \(n=14\).
For an internal index \(j\), put \(q=14-j\) in
\eqref{exapp:prefix-KKT}.  If \(j<8\), the first term is
\(T[f_{14}(q)-g_{14}]>0\), while the others are \(O(C)\).
If \(j=8\), the first term vanishes exactly and the other seven
terms are the strictly positive middle margins displayed above,
multiplied by \(C\).  If \(9\le j\le13\), the ninth group gives
\(Uf_6(q)>0\), with all other terms \(O(T+C)\).
Since \(U=T^3\), all internal multipliers are positive for large
\(T\).  The exterior multipliers vanish and \(\sigma>\theta\),
proving the asserted exact owner partition.

Only groups \(1,9,14\) have growing clocks in the proper owner;
all other binary factors are automatically saturated.  For the edge
\((i,q)=(1,6)\), keep names \(9,\ldots,14\) until their own
deadlines.  Comparing the two completions with the full configuration
and using \(f_{14}(6)=g_{14}\), both reference costs equal
\[
 C\sum_{r=7}^{13}\{\kappa_r(E_r-E_6)-(r-6)\}.
\]
The cuts are late.  The left side makes the first group good, and
both sides make precisely the seven intermediate groups of ranks
\(7,\ldots,13\) good.  In particular the displayed pair gives
\(X_{1,6}\ll\sqrt T\), \(Y_{1,6}\ll1\).

To prove the lower bound for \(X_{1,6}\), consider any legal
completion.  If between one and five names survive through the
rank-six group, its cost is \(Uf_6(q')\gg U\), and all other
possibly negative contributions have total magnitude \(O(T+C)\).
If none survives, the rank-one compensating length \(W\) contributes
nothing, while the first-group contribution is \(g_{14}T\).
If all six survive but name \(14\) ends before its deadline
\(V_{14}\), the rank-one negative contribution is again absent,
leaving \(g_{14}T+O(C)\).  In the remaining case name \(14\)
reaches \(V_{14}\); the completion is late and the first group is
good on the left.  Its score is then at least \(\sqrt T\).
These cases exhaust the templates and prove
\((X_{1,6},Y_{1,6})\asymp(\sqrt T,1)\).

For \((i,q)=(9,0)\), all six deep names end at \(V_8\) or \(V_9\).
At least these six names survive the first group.  Since the number
there lies between \(6\) and \(14\), strict concavity and equal
endpoint values give a contribution at least \(g_{14}T\).  In
group \(9\), the left and right contributions are respectively
\(f_6(0)U=f_6(6)U=0\).  No name survives into the group of length
\(W\).  The remaining costs are \(O(C)\).  These bounds apply to
every completion.  Conversely, ending every earlier name at zero
gives costs \(O(T+C)\), and the good-group correction is at most
\(\sqrt T+11\sqrt C\): the long rank-six group has either no
survivor or all its names.  The exit reserve is bounded.
Thus both endpoint scores have order \(T\).

It remains to exclude additional unsaturated edges on the three
growing clocks.  In group \(1\), the edge \(q=0\) has a bounded
left score and right reference cost \(g_{14}T\), so it is saturated.
For \(1\le q\le5\), the left completion has at most five survivors
in group \(9\).  A nonempty surviving set there costs \(\gg U\);
an empty set leaves \(Tf_{14}(q)+O(C)\gg T\), with no rank-one
compensation.  For \(7\le q\le13\), the entire possible negative
rank-one compensation is \(g_1W=-g_{14}T+O(C)\).
The first-group cost still leaves
\(T[f_{14}(q)-g_{14}]+O(C)\gg T\), and the rank-six contribution
is nonnegative.  In group \(9\), each edge with \(q\ge1\) has
left cost \(Uf_6(q)+O(T+C)\gg U\).  Finally, group \(14\) has
only \(q=0\).  Its left completion keeps name \(14\) through
group \(9\), but supplies no negative rank-one contribution.
If fewer than six names survive group \(9\), its cost is
\(\gg U\); otherwise the first group contributes at least
\(g_{14}T\), up to bounded errors.  Since \(W\asymp T\), this
also saturates the edge.  These cases cover every canonical edge.

The proper native length is \(\asymp U\), and its exit gap is a
fixed positive constant, so its count is \(U^{-\sigma/2}\).
The terminal singleton has length \(R\asymp U\) and fixed positive
drift coefficient \(2^\theta\log2-1\).  Its factor is
\(U^{-1/2}\).  Inserting the two endpoint pairs in
\eqref{profile:edge-factor} proves \eqref{ex:same-answer}.
Direct substitution into \eqref{profile:saddle} gives
\eqref{ex:same-rate}.

\subsection{Two comparisons in different owners}

Write \(z=e^{c/s}\), \(Z=(1+z)^s\), and
\(A=\log(1+z)-z\log z/(1+z)\).  The rank-one equation is
\(F(s;c)=ZA-1\).  At fixed \(c>0\), differentiation gives
\begin{equation}
 \partial_sF(s;c)=Z\left[A^2+
               \frac{c^2}{s^3}\frac{z}{(1+z)^2}\right]>0.
 \label{exapp:rankone-monotone}
\end{equation}
Moreover \(F(1;c)>0\), because
\((1+z)\log(1+z)-z\log z>1\) when \(z>1\).
The rational inequalities
\[
 t_1(17)-\alpha\in(-0.00026097,-0.00026095),\qquad
 t_1(18)-\alpha\in(0.00374983,0.00374984)
\]
give \(F(\alpha;\alpha\log17)<0\).  The roots of
\eqref{ex:multi-roots} are isolated by
\begin{equation}
 0.531542356069<v<0.531542356071,\qquad
 0.5343537358<u<0.5343537362.
 \label{exapp:multi-root-bounds}
\end{equation}
For the first interval substitute its endpoints in
\(F(s;\alpha\log17)\); for the second use the entire first
interval in \(G_2^*\), and substitute the second pair of endpoints
in \(F(s;G_2^*)\).  The signs are respectively negative and positive
under the rational bounds of \S\ref{exapp:certificate-method}.
Equation \eqref{exapp:rankone-monotone} proves uniqueness, and the
enclosures prove the two strict gaps.

On the proposed face, set
\(s_1=u_T\), \(s_2=v_T\), \(s_3=\cdots=s_{18}=\tau_T\).
Solve
\begin{equation}
 \partial_{s_1}J=0,\qquad \partial_{s_2}J=0,\qquad
 \sum_{j=3}^{18}\partial_{s_j}J=0.
 \label{exapp:multi-stationary}
\end{equation}
After division by \(T,T^3,T^{11}\), respectively, the limiting
system is \eqref{ex:multi-roots} together with
\(2^\alpha\log2=1\).  Its Jacobian is triangular when the terminal
variable is solved first.  Its diagonal entries are nonzero by
\eqref{exapp:rankone-monotone} and
\(\frac{d}{d\tau}(2^\tau\log2)|_{\tau=\alpha}=\log2\).
The normalized perturbation in the second equation is \(O(T^{-2})\),
and in the last it is \(O(T^{-6})\); the first is an exact native
equation with perturbed child.  Applying the implicit function
theorem successively gives
\begin{equation}
 u_T-u=O(T^{-2}),\quad v_T-v=O(T^{-2}),\quad
 \tau_T-\alpha=O(T^{-6}).
 \label{exapp:multi-convergence}
\end{equation}

We verify all internal terminal KKT inequalities.  At the limiting
products their leading coefficients, after division by \(T^5\), are
\begin{equation}
 \Lambda_j=\sum_{i=3}^j2^{i-3}
 \left[h_i^{\alpha-1}
   \{h_i\log h_i-(19-j)\log(19-j)\}-(j-i+1)\right],
 \quad 3\le j\le17.
 \label{exapp:multi-KKT}
\end{equation}
Rational lower bounds for these fifteen expressions are in the middle
column below.  The old groups contribute \(O(T^3)\), and parameter
perturbations give smaller errors; hence
\(\lambda_j/T^5=\Lambda_j+O(T^{-2})>0\).  The first two equations
of \eqref{exapp:multi-stationary} give \(\lambda_1=\lambda_2=0\),
and its last equation gives \(\lambda_{18}=0\).  This verifies the
entire KKT system and the exact owner partition.

The same table records the finite signs needed to saturate all
terminal binary edges.  For \(3\le h\le17\), put
\begin{equation}
 \delta_{h,q}=h-1-q-
 h^{\alpha-1}\{h\log h-(q+1)\log(q+1)\},
 \qquad 1\le q\le h-2.
 \label{exapp:terminal-edge-difference}
\end{equation}
The last column bounds \(\min_{1\le q\le h-2}|\delta_{h,q}|\)
from below.  Each entry is obtained by the finite sum
\eqref{exapp:multi-KKT} or the \(h-2\) values in
\eqref{exapp:terminal-edge-difference}, with the rational series
bounds already specified.
\begin{center}
\begin{tabular}{rr|rr}
\(j\)&lower bound for \(\Lambda_j\)&\(h\)&lower bound for
\(\min_q|\delta_{h,q}|\)\\ \hline
3&0.000739&3&0.13787\\
4&0.010639&4&0.16404\\
5&0.068703&5&0.11999\\
6&0.299317&6&0.02515\\
7&1.046964&7&0.10933\\
8&3.195478&8&0.00702\\
9&8.902046&9&0.05276\\
10&23.236999&10&0.06103\\
11&57.730006&11&0.04796\\
12&137.702459&12&0.02228\\
13&316.297189&13&0.01094\\
14&696.898091&14&0.01397\\
15&1448.473445&15&0.01498\\
16&2693.099320&16&0.00755\\
17&3533.706046&17&0.000739
\end{tabular}
\end{center}
In particular all these lower bounds exceed \(7/10000\).

For the singleton owner \([1]\), there is just one paired template,
and no good group.  Its left reference cost is zero.  Its right
reference cost is also zero by \(\partial_{s_1}J=0\).
The left cut is early and the right cut is late.  Since its exit gap
has a positive limit,
\(X_1=1+\widehat R_1\asymp1\), \(Y_1=1\).

The singleton owner \([2]\) also has just one pair, with cuts
\(V_1\) and \(V_2\).  Let \(A_2\) be its native entropy
increment, so that \(A_2Z_2\to1\).  Its left reference cost is
\begin{equation}
 \overline C_{2,-}=T(1-A_2Z_1\beta_1),\qquad
 1-A_2Z_1\beta_1\longrightarrow
          1-(\beta_1^*)^{1-u}>0.
 \label{exapp:old-singleton-score}
\end{equation}
Here the identity
\(Z_1\beta_1=Z_2\beta_1^{1-s_1}\) follows directly from the
recursion, and accounts for the transported old intensity.
The right reference cost is zero by \(\partial_{s_2}J=0\).
The left cut is early since \(T\ll T^3\), and the right cut is
late.  No good group is possible in rank one.  Thus
\(X_2\asymp T\) and \(Y_2=1\).
The two binary ratios are \(\asymp T^{-1}\) and
\(\asymp T^{-2}\), respectively.  Their powers tend to
\(2\chi_1,2\chi_2\).  Local Lipschitz continuity and
\eqref{exapp:multi-convergence} allow replacement by those limiting
powers at a bounded multiplicative cost.  The proper count factors
are \(T^{-u/2}\) and \(T^{-3v/2}\).

For any terminal edge in group \(i\), changing from its left to
its right completion changes cuts only across that group.  Therefore
the reference-cost difference is exactly
\(L_i\delta_{h_i,q}(\tau_T)\), where the notation in
\eqref{exapp:terminal-edge-difference} now uses \(\tau_T\) instead
of \(\alpha\).  The finite lower bounds and
\eqref{exapp:multi-convergence} keep its absolute value
\(\gg L_i\).  If the difference is positive, every right reference
cost is \(\gg L_i\), because every left reference cost is
nonnegative.  If it is negative, every left cost is \(\gg L_i\).
Thus one entire side has this lower bound before taking the two
separate minima.  All \(\sum_{h=3}^{17}(h-2)=120\) binary factors
are saturated; no enumeration of their many individual contexts
is required.

The ordinary suffix starts in group \(18\), since
\(\tau_T\to\alpha_2\).  The common terminal balance gives
\begin{equation}
 D=(2^{\tau_T}\log2-1)T^{11}
   =d_0T^5+O(T^3)+O(T^{-1}),\qquad
 d_0=\sum_{i=3}^{17}2^{i-3}
             \{h_i-1-h_i^\alpha\log h_i\}>0.
 \label{exapp:multi-ordinary}
\end{equation}
Each summand of \(d_0\) is positive because \(\alpha_h>\alpha_2\)
for \(h>2\).  To obtain the error bound, the old groups have total
length \(O(T^3)\), and changing the product in the nonfinal terminal
groups costs \(O(T^5|\tau_T-\alpha|)=O(T^{-1})\).
Now \(U=M=T^{11}\), \(A=B=0\), so
\(\mathcal B_{\rm ord}\asymp T^{-11/2}T^5/T^{11}=T^{-23/2}\).

Finally, \(J_0(T)\) in \eqref{ex:multi-rate} is precisely the
objective at the limiting products \((u,v,\alpha)\).  Its residual
first derivative is zero in the first direction, \(O(T)\) in the
second, and \(O(T^5)\) in the common terminal direction.
Together with \eqref{exapp:multi-convergence}, the linear change is
\(O(T\,T^{-2}+T^5T^{-6})=O(T^{-1})\).
For the quadratic change, differentiating the recursion on a compact
neighborhood gives diagonal bounds \(O(T)\), \(O(T^3)\),
\(O(T^{11})\).  Mixed derivatives involving the first direction
are \(O(T)\), and the remaining mixed derivative is \(O(T^3)\):
only groups preceding both differentiated nodes can contribute.
Their products with the respective increments are no larger than
\(O(T^{-1})\).  Hence \(J(L)-J_0(T)=O(T^{-1})\).
All coordinates tend to infinity, and the positive-vector comparison
proves \eqref{ex:multi-answer}.

To finish the dimensional statement, suppose a terminal owner has
\(C\) names including its permanent child and limiting product
\(\alpha_2\).  A singleton immediately above it solves
\(F(s;\alpha_2\log C)=0\).  By
\eqref{exapp:rankone-monotone}, its product exceeds \(\alpha_2\)
exactly when \(t_1(C+1)>\alpha_2\).  The one-valley property
of Lemma~\ref{hier:lem:valley} and
\(t_1(17)<\alpha_2<t_1(18)\) show that the first integer value
is \(C=17\).  Two preceding singleton owners require
\(H=19\), or \(k=18\).  The root enclosures above verify that
both strict gaps are realized in this first admissible case.
